% =====================================================================
\chapter{Contexte du stage}
\label{ch:contexte}
% =====================================================================

\section{La structure d'accueil}

Talan est un groupe de conseil et d'ingénierie technologique dont l'activité
couvre le conseil en organisation, l'intégration de systèmes et, de plus en
plus, la \emph{data} et l'intelligence artificielle. Ses interventions
s'adressent à plusieurs secteurs, dont la banque, l'assurance, l'énergie et les
transports.

Ce stage s'est déroulé dans le cadre du \textbf{SummerCamp 2026}, le programme
d'été du groupe. Le format est celui d'un stage projet : un sujet ouvert, une
durée courte — environ un mois de travail effectif, du 14 août au 15 septembre
2026 d'après l'historique du dépôt — et une obligation de résultat tangible.
Contrairement à un stage d'observation, il fallait livrer quelque chose qui
tourne : un service déployé, pas une présentation.

Le sujet proposé portait sur le \textbf{risque climatique appliqué à
l'assurance}, un thème qui n'est pas un exercice académique. La France a
effectivement ouvert une discussion de politique publique sur l'adaptation du
régime d'indemnisation des catastrophes naturelles (CatNat) au changement
climatique. Travailler sur la donnée climatique appliquée au bâti, dans ce
contexte, n'est pas un problème inventé pour un stage : c'est un problème que
des professionnels cherchent activement à outiller.

\section{Le problème posé}

\subsection{Trois sources, aucune jointure}

Le point de départ du projet est un constat simple, mais qui n'avait pas été
exploité. La France dispose d'un écosystème open data remarquable sur le
risque :

\begin{itemize}[leftmargin=1.4em]
  \item \textbf{Géorisques} publie les aléas réglementaires : périmètres de
        plans de prévention des risques (PPR), sites et sols pollués, réseaux
        de transport de matières dangereuses, cavités, retrait-gonflement des
        argiles, etc. La donnée est servie en WFS, donc interrogeable
        géométriquement.
  \item La \textbf{BDNB} (Base de Données Nationale des Bâtiments, CSTB) décrit
        chaque bâtiment : usage, hauteur, emprise, matériaux de construction,
        année, performance énergétique, avec des indicateurs de fiabilité.
  \item \textbf{Copernicus} et \textbf{Open-Meteo} fournissent la dimension
        temporelle : projections climatiques à 2050 et 2100, prévisions
        météorologiques horaires, débits de rivière.
\end{itemize}

Aucune de ces sources ne fait la jointure. Géorisques ne sait pas ce qu'est le
bâtiment : il dit qu'un périmètre existe, pas si \emph{ce} bâtiment est dedans.
La BDNB ne sait pas si le bâtiment est inondable. Copernicus travaille sur une
grille d'environ 10~km, donc ne dit rien d'un bâtiment précis.

\begin{encadre}
\textbf{La jointure aléa $\times$ vulnérabilité est le produit.} Un aléa sans
vulnérabilité est une carte. Une vulnérabilité sans aléa est une fiche
technique. Les deux ensemble, pour une même adresse, avec la provenance de
chaque fait, constituent une information qui n'existe nulle part ailleurs.
\end{encadre}

\subsection{Le contournement actuel}

Un souscripteur qui reçoit un dossier d'assurance habitation ou professionnelle
doit aujourd'hui évaluer l'exposition du bien. La pratique observée — décrite
dans les notes de cadrage du projet et confirmée lors des échanges avec un
professionnel du secteur — consiste à ouvrir Géorisques dans un onglet, la BDNB
dans un autre, et à juger visuellement : la rivière est-elle proche ? le
périmètre enveloppe-t-il la parcelle ? le bâtiment est-il en rez-de-chaussée ?

C'est lent, non reproductible, non traçable, et cela ne passe pas l'épreuve
d'un audit. Sous Solvabilité~II, une décision de souscription doit pouvoir être
justifiée. \og J'ai regardé la carte et ça m'a paru bon \fg{} n'est pas une
justification.

\section{La mission}

La mission confiée peut se formuler ainsi :

\begin{quote}
\itshape
Concevoir un service qui, à partir d'une adresse française, produit une fiche
de risque climatique \textbf{au niveau du bâtiment}, en croisant les aléas
réglementaires et la vulnérabilité du bâti, avec une provenance explicite pour
chaque fait — et le rendre utilisable par un souscripteur dans son travail
quotidien.
\end{quote}

Cette formulation est volontairement large. Elle laisse ouvertes trois
questions qui se sont révélées être les vraies décisions du stage :

\begin{enumerate}[leftmargin=1.6em]
  \item \textbf{Quel est le livrable ?} Une interface graphique, un rapport
        narratif, ou une donnée structurée ? La réponse n'était pas évidente
        au départ.
  \item \textbf{Quel est le niveau de résolution acceptable ?} Faire du
        per-bâtiment partout est-il possible ? Sinon, comment le dire
        honnêtement ?
  \item \textbf{Que fait-on de la synthèse ?} Un assureur attend souvent un
        \og score \fg. Est-ce le bon livrable ?
\end{enumerate}

\section{Démarche et organisation du rapport}

Le stage s'est déroulé en trois mouvements, qui structurent ce rapport.

\begin{description}[leftmargin=0em,style=nextline]
  \item[Phase 1 — Exploration (mi-août 2026).] Construction d'une première
        version fonctionnelle : quatre connecteurs de données, une couche de
        notation, un tableau de bord climatique, des extrusions 3D. Le produit
        \emph{marchait}, mais son intérêt était difficile à défendre : il
        reproduisait, à résolution égale et gratuitement, ce que chaque portail
        source fournissait déjà.
  \item[Phase 2 — Recentrage (25 août 2026).] Une remise à plat complète. Le
        livrable a été redéfini comme un contrat de données, le score composite
        a été supprimé, et une grande partie du code de la phase 1 a été
        \emph{supprimée}. Une constitution d'ingénierie a été écrite pour
        verrouiller ces choix par des tests et des règles non négociables.
  \item[Phase 3 — Industrialisation et mise en scène (fin août à
        mi-septembre).] Durcissement du contrat, ajout d'un moteur de scénarios
        interactifs sur la carte, d'un mode de rendu cinématique, et d'un
        déploiement serverless vérifié en ligne.
\end{description}

Le rapport suit cette logique. Le chapitre~\ref{ch:metier} décrit le métier de
l'assurance et ce que le produit doit — ou ne doit pas — être. Le
chapitre~\ref{ch:existant} analyse l'existant et la décision de recentrage. Les
chapitres~\ref{ch:methodo} à \ref{ch:architecture} traitent de la méthode, des
sources de données et de l'architecture. Les chapitres~\ref{ch:scenarios} et
\ref{ch:vfx} décrivent le moteur de scénarios et le mode VFX. Le
chapitre~\ref{ch:qualite} détaille la stratégie de test. Le
chapitre~\ref{ch:deploiement} couvre le déploiement, le
chapitre~\ref{ch:resultats} les résultats mesurés, et le
chapitre~\ref{ch:bilan} le bilan critique.

\begin{figure}[H]
\centering
\begin{tikzpicture}[
  node distance=8mm,
  phase/.style={draw=typhoonBlue,thick,fill=typhoonBlue!6,rounded corners=3pt,
                text width=3.55cm,align=center,minimum height=2.5cm,inner sep=6pt},
  arrow/.style={-{Latex[length=2.4mm]},thick,typhoonBlue}
]
\node[phase] (p1) {\textbf{\color{typhoonDeep}Phase 1}\\[3pt]
  {\small Exploration}\\[4pt]
  {\footnotesize 4 connecteurs, notation, tableau de bord, 3D}\\[3pt]
  {\scriptsize\color{typhoonGrey}14--24 août}};
\node[phase,right=of p1] (p2) {\textbf{\color{typhoonDeep}Phase 2}\\[3pt]
  {\small Recentrage}\\[4pt]
  {\footnotesize Contrat de données, suppression du score, constitution}\\[3pt]
  {\scriptsize\color{typhoonGrey}25--31 août}};
\node[phase,right=of p2] (p3) {\textbf{\color{typhoonDeep}Phase 3}\\[3pt]
  {\small Industrialisation}\\[4pt]
  {\footnotesize Scénarios TRI, mode VFX, tests, déploiement Vercel}\\[3pt]
  {\scriptsize\color{typhoonGrey}1--15 sept.}};
\draw[arrow] (p1) -- (p2);
\draw[arrow] (p2) -- (p3);
\end{tikzpicture}
\caption{Les trois phases du stage. La phase 2 est une phase de soustraction :
elle retire davantage qu'elle n'ajoute, et c'est ce qui a rendu le produit
défendable.}
\label{fig:phases}
\end{figure}
